\documentclass[11pt]{article}
\usepackage{fontspec}
\usepackage[french]{babel}
\usepackage{amsmath}
\usepackage{amssymb}
\usepackage{amsthm}
\usepackage{mathrsfs}
\usepackage{tikz-cd}
\usepackage{enumitem}
\usetikzlibrary{babel}

\newcommand{\printedpage}[1]{\marginpar{\footnotesize p.~#1}}
\newcommand{\msp}[1]{\marginpar{\footnotesize [#1]}}

\begin{document}

\printedpage{24}
% head: A. Grothendieck : Dérivateurs, Ch. XVII, Édition des Universités de Montpellier II et Paris VII

\msp{45}dans $\mathcal{M}(I^\circ)$, considérons les triples
\[
(J,\ Z,\ i,\ p)\ ,
\]
où $J\subset I$ est un ouvert, $Z$ un objet de $\mathcal{M}(J^\circ)$, et $i$, $p$ sont des flèches
\[
\left\{
\begin{aligned}
& X|J \xrightarrow{\ i\ } Z \xleftarrow{\ p\ } Y|J \\
& \text{avec } i\in C_{J^\circ},\ p\in W_{J^\circ}.
\end{aligned}
\right.
\]
On ordonne les quadruples par prolongement. Si on a un ensemble totalement ordonné, de $(J_\alpha,Z_\alpha,i_\alpha,p_\alpha)$, posant $J=\bigcup I_\alpha$, on trouve de $Z$, $i$, $p$ uniques sur $J$ qui prolongent % [sic] \bigcup I_\alpha -> \bigcup J_\alpha, the union is of the open sets J_\alpha
les $Z_\alpha$, $i_\alpha$, $p_\alpha$ donnés, de plus on aura $i\in C_{J^\circ}$, $p\in W_{J^\circ}$, car l'appartenance à $C_{J^\circ}$ resp. $W_{J^\circ}$ est question locale sur $J$. Donc par Zorn, on est ramené à prouver que si $(J,Z,i,p)$ est maximal, on a $J=I$ – d'ou la factorisation voulue $f=pi$. Donc à prouver que si % [sic] d'ou -> d'où, missing accent
$J\neq I$, le quadruple n'est pas maximal. Considérons le fermé non vide $I\backslash J$. $I$ est supposé artinien, il y a un \textit{élément minimal} $i_0$ de $I\backslash J$. Posons
\msp{46}
\[
J' \;=\; J\cup\{i_0\} \;=\; J\cup I/i_0,
\]
et montrons qu'on peut prolonger le triplet $Z$, $i$, $p$ en $Z'$, $j'$, $p'$ sur $J'$. Posant $K'=I/i_0$, % [sic] j' -> i', the triplet is Z, i, p
$K=J\cap K=I_{<i_0}$, on a donc % [sic] J\cap K -> J\cap K', as in the diagram below
\[
\begin{tikzcd}
K=J\cap K' \arrow[r, hook] \arrow[d, hook] & K'=I/i_0 \arrow[d, hook] \\
J \arrow[r, hook] & J'=J\cup K'.
\end{tikzcd}
\]

On voit alors que les prolongements de $(Z,i,p)$ de $J$ à $J'$ correspondent à ceux de $(Z|K',i|K',p|K')$ à $K$. Pour ceci, on note que pour une flèche $\varphi$ d'objets de $\mathcal{M}(J'^\circ)$, cette flèche est dans $C_{J'^\circ}$ resp. $W_{J'^\circ}$, si et seulement si ses restrictions à $J$ et à $K'$ le sont.

On est donc ramené au cas standard d'un ensemble ordonné avec objet final $e_0$, ensemble ordonné qu'on va noter $I$, et l'ensemble ordonné $I'$ des objets $i<e_0$, i.e. $I'=I\backslash\{i_0\}$. % [sic] I\backslash\{i_0\} -> I\backslash\{e_0\}, I' consists of the i < e_0
On a $f:X_\bullet\longrightarrow Y_\bullet$, flèche dans $\mathcal{M}(I^\circ)$, une factorisation de $f|I'$ en $p'\circ i'$ avec $i'\in C_{I'^\circ}$, $p'\in W_{I^\circ}$, % [sic] W_{I^\circ} -> W_{I'^\circ}, factorisation of f|I' on I'
\msp{47}
et il s'agit de la prolonger en une factorisation similaire de $f$. On note que la donnée d'un objet $X_\bullet$ de $\mathcal{M}(I^\circ)$ équivaut à celle d'un objet $X'_\bullet$ de $\mathcal{M}(I'^\circ)$ (la restriction de $X$ à $I'$) et d'une flèche
\[
X'_{I'} \xrightarrow{\ \alpha_X\ } X_{e_0}=X_0
\]
dans $\mathcal{M}$, où $X'_{I'}=\varinjlim\limits_{i'\in I'}X'_{i'}$. Donc la donnée de la flèche $f:X_\bullet\longrightarrow Y_\bullet$ revient à celle de $f'=f|I':X'_\bullet\longrightarrow Y'_\bullet$, et d'une flèche $f_0:X_0\longrightarrow Y_0$ rendant commutatif le carré
\printedpage{25}
% head: A. Grothendieck : Dérivateurs, Ch. XVII, Édition des Universités de Montpellier II et Paris VII
\[
\begin{tikzcd}
X'_{I'} \arrow[r, "\alpha_X"] \arrow[d, "f'_{I'}"'] & X_0 \arrow[d, "f_0"] \\
Y'_{I'} \arrow[r, "\alpha_Y"] & Y_0\ .
\end{tikzcd}
\]

De même, trouver une factorisation de $f$ en $pi$, qui prolonge une factorisation donnée $f'=p'i'$ sur $J'$, revient à la donnée d'un diagramme commutatif
\[
\begin{tikzcd}[column sep=small, row sep=small]
& X'_{I'} \arrow[rr, "\alpha_X"] \arrow[dl, "i'_{I'}"'] \arrow[dd, "f'_{I'}"] & & X_0 \arrow[dl, "i_0"'] \arrow[dd, "f_0"] \\
Z_{I'} \arrow[rr, "\alpha_Z"'] \arrow[dr, "p'_{I'}"'] & & Z_0 \arrow[dr, "p_0"'] & \\
& Y'_{I'} \arrow[rr, "\alpha_Y"'] & & Y_0\ ,
\end{tikzcd}
\tag*{$(*)$}
\]
\msp{48 vide}
\msp{49}
où l'objet $Z_0$, et les trois flèches $\alpha_Z$, $i_0$, $p_0$ sont à déterminer. Si on a $i'\in C_{I'^\circ}$, $p\in W_{I'^\circ}$, % [sic] p -> p', the given factorisation is f'=p'i'
alors on conclut par théorème 2.1. $1^\circ$ (cor. 2.2) que
\[
i'_{I'}\ \in\ C\ ,
\]
et si on veut que $i\in C_{I^\circ}$, $p\in W_{I^\circ}$, il reste seulement à exiger
\[
p_0\in W, \qquad (i'_{I',I}:\underbrace{Z_{I'}\amalg_{X_{I'}}X_0}_{Z'_0}\longrightarrow Z_0)\in C,
\]
d'après le théorème 2.1. $2^\circ$. On va alors construire le diagramme $(*)$, à partir du diagramme similaire ci-dessous
\[
\begin{tikzcd}
& X_{I'} \arrow[rr, "\alpha_X"] \arrow[dl, "i'_{I'}\in C"'] \arrow[dd, "\text{cocart.}" {pos=0.25}, "f'_{I'}" {pos=0.75}] & & X_0 \arrow[dl, "i'_0\in C"'] \arrow[dd, "f_0" {pos=0.25}] \arrow[dr, dotted, "i_0=j_0i'_0"] & \\
Z_{I'} \arrow[rr, "\alpha'_Z"] \arrow[dr, "p'_{I'}"'] & & Z'_0=Z_{I',I} \arrow[rr, dotted, "j_0\in C"] \arrow[dr, "p'_0"'] & & Z_0 \arrow[dl, dotted, "p_0\in W"] \\
& Y_{I'} \arrow[rr, "\alpha_Y"] & & Y_0 &
\end{tikzcd}
\]
(avec $Z_0$ de $(*)$ remplacé par $Z'_0=Z_{I',I}$) en factorisant la flèche canonique $p'_0:Z'_0\longrightarrow Y_0$ en
\[
p'_0 \;=\; p_0j_0\ , \qquad \text{avec } p_0\in W,\ j_0\in C\ ,
\]
et en posant
\[
i_0 \;=\; j_0i'_0\ ,
\]
ce qui achève la construction.

On a ainsi prouvé le

\printedpage{26}
% head: A. Grothendieck : Dérivateurs, Ch. XVII, Édition des Universités de Montpellier II et Paris VII

\msp{50}\textbf{Théorème 3.1.} \textit{Soit $(\mathcal{M},W,C)$ une catégorie à cofibrations (cf. page \textup{39}), $\mathcal{M}$ stable par $\varinjlim$ finies. Pour tout ensemble odonné $I$, considérons le tuple $(\mathcal{M}(I^\circ),W_{I^\circ},C_{I^\circ})$ correspondant, où $C_{I^\circ}$ est défini au paragraphe \textup{2} (cf. à ce sujet le théorème \textup{2.1}). Supposons que $I$ soit un ensemble ordonné \textup{artinien} (p.ex. $I$ localement fini, i.e. les $I/i$ finis). Supposons de plus qu'on soit dans l'un des deux cas suivants :} % [sic] odonné -> ordonné, missing letter
\begin{itemize}[leftmargin=2.5em]
\item[a)] \textit{L'espace \textup{Esp}$(I)$ est noethérien, i.e. toute suite croissante d'ouverts de $I$ est stationnaire (p.ex. $I$ fini).}
\item[b)] \textit{$C$ et $TC=C\cap W$ sont stables, non seulement par composition ordinaire, mais aussi par composition transfinie. Dans ce cas, on suppose aussi $\mathcal{M}$ stable par petites $\varinjlim$ filtrantes (donc par petites $\varinjlim$ quelconques).}
\end{itemize}
\textit{Sous ces conditions, $(\mathcal{M}(I^\circ),W_{I^\circ},C_{I^\circ})$ est également une catégorie à cofibrations.}

\textbf{Proposition 3.2.} \textit{Soit $u:I\longrightarrow I'$ une application croissante entre ensembles ordonnés, d'où un couple de foncteurs adjoints}
\[
\begin{tikzcd}
\mathcal{M}(I^\circ) \arrow[r, shift left, "u_{!}"] & \mathcal{M}(I'^\circ) \arrow[l, shift left, "u^*"]
\end{tikzcd}\ .
\tag*{(3.2.1)}
\]
\msp{51}
\begin{itemize}[leftmargin=2.5em]
\item[a)] \textit{Le foncteur $u^*$ est compatible avec les localiseurs $W_{I'^\circ}$, $W_{I^\circ}$. Si $I$ est discrète, ou $u$ étale, alors $u^*$ applique $C_{I'^\circ}$ dans $C_{I^\circ}$. Il en est encore de même quand $u$ est localement cofiltrant, i.e. si pour tout $i'\in\operatorname{Ob}I'$, la catégorie $i'\backslash I$ est localement filtrante, i.e. ses composantes connexes sont cofiltrantes (condition $\mathrm{PS}_{\mathrm{g}}$ \textup{1} sur $i'\backslash I$)} \textsuperscript{(}\footnote[16]{inutile de parler de ``faiblement'' localement cofiltrant, car il s'agit d'ensembles ordonnés où il n'y a pas à distinguer entre notions forte et faible.}\textsuperscript{)}.
\item[b)] \textit{Dans tous les cas, $u_!$ transforme $C_{I^\circ}$ en $C_{I'^\circ}$, $TC_{I^\circ}$ en $TC_{I'^\circ}$, enfin $\mathcal{M}(I^\circ)_{\mathrm{c}}$ en $\mathcal{M}(I'^\circ)_{\mathrm{c}}$.}
\end{itemize}

\textsc{Démonstration}. Prouvons a), dans le cas le plus gênant, où on suppose seulement que $u$ est localement faiblement cofiltrant. On applique II 7.4 (p. 92), qui nous dit que $\forall i\in\operatorname{Ob}I$,
\msp{52}
considérant le morphisme induit par $u$
\[
I/i \longrightarrow I'/i' \qquad (\text{où } i'=u(i)),
\]
celui-ci est \textit{(globalement)} faiblement cofiltrant, et que pour tout ouvert $U$ de $I/i$, désignant par $U'$ l'ouvert engendré par son image, le morphisme induit $v:U\longrightarrow U'$ est encore faiblement cofiltrant, et a fortiori \textit{cofinal.} Soit alors $f':X'\longrightarrow Y'$ une flèche dans $C_{I'^\circ}$, prouvons que
\[
\underbrace{u^*(f')}_{f}\ :\ \underbrace{u^*(X')}_{X}\ \longrightarrow\ \underbrace{u^*(Y')}_{Y}
\]
est dans $C_{I^\circ}$, i.e. que pour tout $i\in I$, et $U$ ouvert de $I/i$, dans le diagramme canonique
\printedpage{27}
% head: A. Grothendieck : Dérivateurs, Ch. XVII, Édition des Universités de Montpellier II et Paris VII
\[
\begin{tikzcd}
X_U \arrow[r] \arrow[d] \arrow[dr, phantom, "\text{cocart.}"] & X(i) \arrow[d] \\
Y_U \arrow[r] & Y_{U,i} \arrow[d, "v_{U,i}"] \\
& Y(i)\ ,
\end{tikzcd}
\tag*{$(*)$}
\]
\msp{53}
la flèche $v_{U,i}$ est dans $C$. Le diagramme $(*)$ s'envoie dans le diagramme similaire pour $f':X'\longrightarrow Y'$, $i'=u(i)$, et l'ouvert $U'\subset I'/i'$ engendré par $u(U)$ :
\[
\begin{tikzcd}[column sep=small, row sep=small]
X_U \arrow[rr] \arrow[dr, "f_U"] \arrow[dd, "c_U^{X'}"'] & & X(i) \arrow[dr] \arrow[dd, "\text{cocart.}"' {pos=0.2}, "\mathrm{id}"' {pos=0.75}] & & \\
& Y_U \arrow[rr] \arrow[dd, "c_U^{Y'}"' {pos=0.3}] & & Y_{U,i} \arrow[dr, "v_{U,i}"] \arrow[dd, "c^{f'}_{U,i}"] & \\
X'_{U'} \arrow[rr, "\text{cocart.}"' {pos=0.6}] \arrow[dr, "f'_{U'}"'] & & X'(i') \arrow[dr] & & Y(i) \arrow[dd, "\mathrm{id}"] \\
& Y'_{U'} \arrow[rr] & & Y'_{U',i'} \arrow[dr, "v_{U',i'}"'] & \\
& & & & Y'(i')\ .
\end{tikzcd}
\tag*{$(**)$}
\]

Par cofinalité, les deux flèches verticales $c_U^{X'}$ et $c_U^{Y'}$ sont des isomorphismes, en plus de $X(i)\xrightarrow{\ id\ }X'(i')$, donc la quatrième flèche du cube, $c^{f'}_{U,i}$, est un isomorphisme. Comme $v_{U',i'}$ est dans $C$, $v_{U,i}$ aussi, on gagne.

Prouvons b). Soit donc
\[
f\ :\ X\longrightarrow Y
\]
une flèche dans $C_I$, et soit
\[
f' \;=\; u_!(f)\ :\ u_!(X)=X' \longrightarrow u_!(Y)=Y',
\]
prouvons que c'est dans $C_{I'^\circ}$, donc que
\msp{54}
pour deux ouverts $U'\subset V'$ de $I'$ (où on peut supposer que $V'=I'/i'$), dans le diagramme-base de $(**)$, $v_{U',V'}$ soit dans $C$. Soit $U=u^{-1}(U')$, $V=u^{-1}(V')$, et considérons le % ?? uncertain: hyphen at line end in "diagramme-base" -- kept as a compound hyphen, or "diagrammebase"
diagramme similaire à $(**)$, défini par les deux inclusions d'ouverts $U'\subset V'$, $U\subset V$. Cette fois, comme la formation de $u_!$ commute aux restrictions aux ouverts, par transitivité des $\varinjlim$ , les trois flèches vertivales $c^X_{U'}$, $c^Y_{U'}$, $c^X_{V'}$ sont des isomorphismes, donc aussi $c^f_{U',V'}$. De % [sic] vertivales -> verticales, misprint
même, la flèche $Y_V\longrightarrow Y'_{V'}$ (à la place de $Y(i)\longrightarrow Y'(i')$) est un isomorphisme. Comme $v_{U,V}\in C$, il en est de même de $v_{U',V'}$, ce qui prouve que $f'\in C_{I'^\circ}$. On procède de même pour $TC_I$. Enfin, l'assertion que $u_!$ applique $\mathcal{M}(I^\circ)_{\mathrm{c}}$ dans $\mathcal{M}(I'^\circ)_{\mathrm{c}}$ résulte des définitions et de $u_!(C_{I^\circ})\subset C_{I'^\circ}$.

\textbf{Corollaire 3.3} (de prop. 3.2 et th. 3.1). \textit{Soit $u:I\longrightarrow I'$ comme dans \textup{3.2}. Supposons de plus $I$ artinien, et}
\printedpage{28}
% head: A. Grothendieck : Dérivateurs, Ch. XVII, Édition des Universités de Montpellier II et Paris VII
\msp{55}
\textit{qu'on soit pour $\mathcal{M}$, $W$, $C$, $I$ dans l'une des conditions \textup{a), b)} du théorème \textup{3.1}. Alors $u_!$ induit une flèche}
\[
u_{!\,\mathrm{cof}}\ :\ \mathcal{M}(I^\circ)_{\mathrm{c}} \longrightarrow \mathcal{M}(I'^\circ)_{\mathrm{c}}\ \subset\ \mathcal{M}(I'^\circ)
\]
\textit{qui transforme flèches de $W(I^\circ)$ en flèches de $W(I'^\circ)$ (i.e. compatible avec les localiseurs). Il induit donc un foncteur sur les localisés, d'où un}
\[
u_!^W\ :\ \mathrm{Ho}_W(I^\circ) \longrightarrow \mathrm{Ho}_W(I'^\circ) \qquad \mbox{\textsuperscript{(}\footnotemark[17]\textsuperscript{)}}.
\tag*{(3.3.1)}
\]
\footnotetext[17]{\textbf{N.B.} Cette flèche est définie, même sans supposer $I'$ ordonné, car $u_{!\,\mathrm{cof}}$ transforme encore flèches de $W(I^\circ)$ en flèches de $W(I'^\circ)$. Cf. commentaire p. 86.}

\textsc{Démonstration}. Par le théorème 3.1, on sait que $\mathcal{M}(I^\circ)$, $W(I^\circ)$, $C(I^\circ)$ est une structure de catégorie de modèles à cofibrations. On sait par ailleurs que $TC_{I^\circ}=C_{I^\circ}\cap W(I^\circ)$, par th. 2.1. $3^\circ$ (p. 27). Donc par 3.2, $u_!$ applique $C_{I^\circ}\cap W(I^\circ)$ dans $TC_{I'^\circ}=C_{I'^\circ}\cap W(I'^\circ)$, et a fortiori, en des flèches dans $W(I'^\circ)$. Par ailleurs, par la théorie des catégories à cofibrations, on trouve que toute flèche $f$ entre objets cofibrants est composée $f=pi$, où $i\in C_{I^\circ}$, et où $p$ est dans $W(I^\circ)$, et de plus \textit{inverse à gauche} d'une flèche dans $C_{I^\circ}\cap W(I^\circ)$, i.e. d'une flèche de $(TC)_{I^\circ}$. Quand
\msp{56}
$f\in W(I^\circ)$, alors $i\in W(I^\circ)$, donc $i\in W(I^\circ)\cap C_{I^\circ}=(TC)_{I^\circ}$. Au total, toute flèche $f\in W(I^\circ)_{\mathrm{c}}$ est composée $pi$ dans $\mathcal{M}(I^\circ)_{\mathrm{c}}$, où $i\in(TC)_{I^\circ}$, et $p$ est inverse à gauche d'une $j\in(TC)_{I^\circ}$. Comme $u_!$ transforme $(TC)_{I^\circ}$ en $(TC)_{I'^\circ}$, il transforme $f$ en $f'=p'i'$, avec $i'\in(TC)_{I'^\circ}\subset W_{I'^\circ}$, et $p'$ inverse à gauche d'une $j'=u_!(j)$ qui est également dans $W_{I'^\circ}$. Mais alors $i'\in W(I'^\circ)$, donc aussi $f'=p'i'\in W_{I'^\circ}$, ce qui prouve que $u_!(W(I^\circ)_{\mathrm{c}})\subset W(I'^\circ)$. La dernière assertion provient du lemme suivant, appliqué à $(\mathcal{M}(I^\circ),W(I^\circ),C_{I^\circ})$ au lieu de $(\mathcal{M},W,C)$ :

\textbf{Lemme 3.4.} \textit{Soit $(\mathcal{M},W,C)$ une catégorie à cofibrations. Alors l'inclusion $\mathcal{M}_{\mathrm{c}}\subset\mathcal{M}$ induit une équivalence}
\[
\mathcal{M}_{\mathrm{c}}W_{\mathrm{c}}^{-1} \xrightarrow{\ \sim\ } \mathcal{M}W^{-1}=\mathrm{Ho}_W\ .
\]

Nous donnerons la démonstration (par AQT \texttt{[âne qui trotte]}) ultérieurement.

Cela achève la démonstration de 3.3.

\textbf{3.5.} Il se pose la question si le foncteur $u_!^W$ de (3.3.1) est bien un adjoint à gauche
\msp{57}
du foncteur en sens inverse $u^*_W$, induit par $u^*$. Notons que c'est le cas si $u^*$ transforme $\mathcal{M}(I'^\circ)_{\mathrm{c}}$ en $\mathcal{M}(I^\circ)_{\mathrm{c}}$, et quand on suppose que $I'$ est également artinien et satisfait également à l'une des hypothèses a) b) du théorème 3.1. Car alors $(\mathcal{M}(I^\circ),W(I^\circ),C_{I^\circ})$ est une catégorie à cofibrations par le théorème 3.1, on peut lui appliquer le lemme précédent, donc $u^*_W$ et $u_!^W$ se réalisent à partir d'un couple de foncteurs
\[
\begin{tikzcd}
\mathcal{M}(I^\circ)_{\mathrm{c}} \arrow[r, shift left, "u_{!}^{\mathrm{c}}"] & \mathcal{M}(I'^\circ)_{\mathrm{c}} \arrow[l, shift left, "u_{\mathrm{c}}^*"]
\end{tikzcd}\ ,
\]

\end{document}
